\section{Introduction}
\label{sec:intro}

A world model is supposed to discard what does not matter. Joint-embedding predictive architectures
(JEPAs) make this a learning problem: an encoder and a predictor are trained together so that the
predictor forecasts the encoder's own next output, and nothing in the objective asks for pixels
back \citep{lecun2022path,assran2023ijepa,bardes2024vjepa,maes2026lewm}. The promise for physical
intelligence is a compact, reusable description of the world, free of texture and lighting. The
difficulty is that the same objective decides two things at once: \emph{what} to encode and
\emph{how} to predict it. We show that in stochastic worlds it can decide both against the
controller, for two reasons that are independent of each other and fail in different places.

The first reason is representational. Squared next-embedding error is minimized by keeping
predictable features and dropping unpredictable ones. A background fixed for an episode is
perfectly predictable even if the object that the agent controls is ignored
\citep{sobal2022slow}. Less obviously, the controlled object itself becomes the unpredictable part
as soon as its dynamics carry fresh innovations that no history reveals. The objective then has a
reason to \emph{learn to ignore what matters}, and it does so without collapsing: the latent stays
non-degenerate, the Gaussian regularizer is satisfied and the prediction loss is low.

The second reason is predictive. A squared loss produces the conditional mean of the next latent.
Two exact calculations show what a planner loses by working with that mean even when it is correct.
An action whose outcomes straddle the goal has its mean exactly on the goal, so a planner scoring
the mean prefers it to an action that reliably lands close. And an exact one-step mean, applied
recursively, need not equal the multi-step mean whenever the dynamics are nonlinear. Neither
failure is an optimization failure, and no better estimate of the same mean repairs it.

We make both mechanisms measurable in a deliberately small visual world: one red cube, four
movement actions, a background that is either redrawn every frame or held for the episode, and
optional stochastic wind or an approaching wall whose opening is revealed late. Position, goal and
executed path are directly visible, so representation and planning can be separated without
leaning on an aggregate game score. The main empirical findings are:
\begin{enumerate}
\item \textbf{Nuisance persistence reverses under stochastic dynamics.} With deterministic motion,
per-frame background noise gives the more accurate position read-out ($0.115$ versus $0.764$
pixels). With hidden wind the ordering flips and the per-frame model loses the cube altogether
($21.33$ versus $0.668$ pixels). Architecture, latent dimension, data, updates and three model
seeds are matched. A pixel policy trained inside the failing model's imagination dies in every
evaluation episode, sooner than a blind random walk, because the model's own termination head
reads death on every latent its predictor produces; the same policy inside a flow model trained on
the same data survives $0.13$ to $0.95$ of the episodes (Section~\ref{sec:representation}).
\item \textbf{A latent that holds the state can still fail imagination.} In the gate world the
frozen JEPA latent, read through its own decoder, supports survival when the policy is trained on
real transitions; policies trained inside the JEPA's imagination with the same network, budget and
evaluation keys \pendingresult{state the shared-CNN outcome}. A flow-matching model trained on the
same data is the distributional reference, and decoded-real controls qualify both interfaces before
any imagination result is read (Section~\ref{sec:planning}).
\item \textbf{The consequence survives outside the toy.} In Craftax, under matched data and
training budgets and a shared pixel policy, \pendingresult{state the inventory-campaign outcome}
(Section~\ref{sec:craftax}).
\end{enumerate}
Section~\ref{sec:mechanisms} gives the two mechanisms with their exact examples and the Gaussian
nuisance construction that makes the first one precise. The conclusion is not that JEPAs cannot
model stochastic worlds; it is that avoiding representation collapse and predicting latents
accurately are insufficient for stochastic control, and that fixing the predictor helps only if the
representation has kept the state in the first place.
